\documentclass[border=10pt]{standalone}
\usepackage[no-math]{fontspec}
\usepackage[brazil]{babel}
\usepackage{amsmath,amssymb}
\IfFileExists{C:/Windows/Fonts/Lato_Regular.ttf}{%
  \setmainfont{Lato}[
    Path = C:/Windows/Fonts/,
    UprightFont = Lato_Regular.ttf,
    BoldFont = Lato_Bold.ttf,
    ItalicFont = Lato_Italic.ttf,
    BoldItalicFont = Lato_BoldItalic.ttf
  ]
  \setsansfont{Lato}[
    Path = C:/Windows/Fonts/,
    UprightFont = Lato_Regular.ttf,
    BoldFont = Lato_Bold.ttf,
    ItalicFont = Lato_Italic.ttf,
    BoldItalicFont = Lato_BoldItalic.ttf
  ]
}{%
  \RequirePackage[default]{lato}
}
\renewcommand{\familydefault}{\sfdefault}

\usepackage{xcolor}
\usepackage{tikz}
\usepackage{pgfplots}
\pgfplotsset{compat=1.18}
\usetikzlibrary{arrows.meta,positioning,shapes.geometric,fit,backgrounds,calc,decorations.pathreplacing,patterns}

% Paleta institucional UEFS / Sociologia Rural
\definecolor{srNavy}{HTML}{1F4268}
\definecolor{srNavyDark}{HTML}{152E4A}
\definecolor{srGold}{HTML}{B38A3D}
\definecolor{srGoldDark}{HTML}{8C6927}
\definecolor{srTeal}{HTML}{0E7C7B}
\definecolor{srForest}{HTML}{2D6A4F}
\definecolor{srAlert}{HTML}{9E2A2B}
\definecolor{srInk}{HTML}{26313D}
\definecolor{srMuted}{HTML}{5B6573}

\definecolor{bgNavy}{HTML}{EDF2F7}
\definecolor{bgGold}{HTML}{F8F3E9}
\definecolor{bgTeal}{HTML}{EBF6F5}
\definecolor{bgForest}{HTML}{EEF5F0}
\definecolor{bgAlert}{HTML}{FDF2F2}

\tikzset{
  >=Stealth,
  srBox/.style={
    draw=#1, fill=white, rounded corners=5pt, line width=1.15pt,
    inner sep=7pt, align=center, font=\footnotesize\sffamily, text=srInk
  },
  srTint/.style={
    draw=#1, fill=#1!8!white, rounded corners=5pt, line width=1.15pt,
    inner sep=7pt, align=center, font=\footnotesize\sffamily, text=srInk
  },
  srBadge/.style={
    fill=#1, text=white, rounded corners=9pt, inner xsep=8pt, inner ysep=3pt,
    font=\scriptsize\bfseries\sffamily
  },
  srTitle/.style={font=\normalsize\bfseries\sffamily, text=srNavy, align=center},
  srSub/.style={font=\footnotesize\sffamily, text=srMuted, align=center},
  srArrow/.style={->, >=Stealth, line width=1.35pt},
  srDArrow/.style={<->, >=Stealth, line width=1.35pt}
}


\begin{document}
\begin{tikzpicture}
  % Título principal
  \node[srTitle, font=\large\bfseries\sffamily] (title) at (0, 3.65)
    {Paradigmas de Desenvolvimento Rural: Do Produtivismo ao Ecodesenvolvimento};
  \node[srSub] (sub) at (0, 3.12)
    {Ruptura analítica entre o modelo técnico-produtivista e a abordagem multidimensional do território};

  % Dimensões das colunas
  \def\colW{6.5cm}
  \def\colH{5.25cm}

  % Coluna Esquerda: Paradigma Clássico / Revolução Verde
  \node[draw=srAlert, fill=bgAlert, rounded corners=6pt, line width=1.25pt,
        minimum width=\colW, minimum height=\colH, anchor=north] (leftCard) at (-4.35, 2.55) {};
  \node[fill=srAlert, text=white, rounded corners=4pt, text width=5.9cm,
        inner ysep=6pt, font=\small\bfseries\sffamily, align=center, anchor=north]
        (leftHdr) at ($(leftCard.north)+(0,-0.16)$)
    {Paradigma Clássico / Revolução Verde};
  \node[text=srInk, font=\footnotesize\sffamily, text width=5.85cm, align=left, anchor=north]
        at ($(leftHdr.south)+(0,-0.22)$) {%
    \textbf{Lógica:} modernização conservadora e pacote tecnológico exógeno.\par\vspace{7pt}
    $\bullet$\;\;Foco exclusivo em PIB agrícola e produtividade por hectare\par\vspace{6pt}
    $\bullet$\;\;Mecanização pesada, monocultura e pacote químico (agrotóxicos e fertilizantes)\par\vspace{6pt}
    $\bullet$\;\;Concentração fundiária, endividamento e êxodo rural acelerado
  };

  % Coluna Direita: Desenvolvimento Territorial Sustentável
  \node[draw=srForest, fill=bgForest, rounded corners=6pt, line width=1.25pt,
        minimum width=\colW, minimum height=\colH, anchor=north] (rightCard) at (4.35, 2.55) {};
  \node[fill=srForest, text=white, rounded corners=4pt, text width=5.9cm,
        inner ysep=6pt, font=\small\bfseries\sffamily, align=center, anchor=north]
        (rightHdr) at ($(rightCard.north)+(0,-0.16)$)
    {Desenvolvimento Territorial Sustentável\\[-1pt]{\scriptsize\normalfont\sffamily (Ignacy Sachs; Amartya Sen)}};
  \node[text=srInk, font=\footnotesize\sffamily, text width=5.85cm, align=left, anchor=north]
        at ($(rightHdr.south)+(0,-0.22)$) {%
    \textbf{Lógica:} sustentabilidade multidimensional e ampliação de capacidades.\par\vspace{7pt}
    $\bullet$\;\;5 Dimensões (Sachs): Social, Econômica, Ecológica, Espacial e Cultural\par\vspace{6pt}
    $\bullet$\;\;Expansão de liberdades substantivas, autonomia e equidade (Sen)\par\vspace{6pt}
    $\bullet$\;\;Valorização de mercados locais, agroecologia e diversidade sociocultural
  };

  % Seta de transição entre os paradigmas
  \draw[srArrow, color=srGoldDark, line width=1.8pt]
    ($(leftCard.east)+(0.18,0)$) -- ($(rightCard.west)+(-0.18,0)$);
  \node[fill=srGoldDark, text=white, rounded corners=4pt, inner xsep=8pt, inner ysep=3.5pt,
        font=\scriptsize\bfseries\sffamily, anchor=south] at ($(leftCard.east)!0.5!(rightCard.west)+(0,0.22)$) {Transição};
\end{tikzpicture}
\end{document}
